% !TeX root = ../../Book1.tex
\section{特殊有界变差函数与紧性}
\label{b1:sec:2804}

在允许界面出现的变分问题中，变化的代价通常分成两类：
区域内部的梯度能量，以及界面本身的面积代价。
如果把所有有界变差函数都作为候选，而能量只计入这两项，
Cantor 部分便可能不付出任何代价。
因此需要一个排除该部分、同时保留跳跃的函数类。

\ProseParagraphBreak

本节先定义特殊有界变差函数，说明它对有界变差空间中的严格收敛并不闭合，
再记录带有值域、超线性梯度和跳跃面积控制的 Ambrosio 紧性定理。
定理用到\chapref{b1:ch:27}中有界变差函数的紧性，以及前节作为深层输入记录的 Volpert 链律。
后者如何用于特殊有界变差函数空间的闭合性，可参看 Alberti、Mantegazza 的
《A Note on the Theory of SBV Functions》%
\cite[Theorem 1.4, Propositions 2.3 and 2.5]{AlbertiMantegazza1997SBV}。
% 实读作者公开版第4--7页，对应刊印378--381页；
% Theorem1.4只陈述紧性及绝对连续/跳跃部分分开收敛，不包含面积下半连续性。
% 本节取其增长函数f(t)=1(t>0),f(0)=0，采用正弦/余弦测试展开闭合桥。

\subsection{特殊有界变差函数}
\label{b1:subsec:280401}

\begin{definition}\label{b1:def:sbv-function}
设 \(u\in BV(\Omega)\)。若 \(D^cu=0\)，则称 \(u\) 为%
\term[teshu-youjie-biancha-hanshu]{特殊有界变差函数}[special function of bounded variation]，
记 \(u\in SBV(\Omega)\)。\foreignidx{SBV}因此
\[
 Du=\nabla u\,m_d+(u^+-u^-)\nu_u\,
 \mathcal H^{d-1}\mathbin{\llcorner}J_u.
\]
\end{definition}
\symidx{\(SBV(\Omega)\)：特殊有界变差函数空间}

每个 \(W^{1,1}\) 函数都属于 \(SBV\)。
有界区域内有限周长集合的示性函数也属于 \(SBV\)：
它的导数由约化边界上的面积表示，故完全属于跳跃部分。
但特殊有界变差函数空间既不等于 Sobolev 空间，也不等于整个有界变差函数空间：
一个非零阶跃不属于 \(W^{1,1}\)，而前节的连续 Cantor 函数不属于 \(SBV\)。
定义本身也不要求跳跃集具有有限面积，更不要求它是闭集或由有限张光滑曲面组成。

\begin{proposition}\label{b1:prop:sbv-not-strictly-closed}
在区间 \((0,1)\) 上，特殊有界变差函数空间对有界变差空间中的严格收敛并不闭合。
甚至存在 \(0\le u_n\le1\) 的连续分段线性函数，
满足 \(u_n\to F_C\) 一致收敛、\(|Du_n|\bigl((0,1)\bigr)=1\)，
而 \(D^cF_C\ne0\)。
\end{proposition}
\begin{proof}
记 \(C_n\) 为三分 Cantor 集构造中的第 \(n\) 层闭区间之并。
在每个长度 \(3^{-n}\) 的基本区间上，用直线连接 Cantor 函数的两个端点值；
在其余间隙上保持 Cantor 函数的常值，所得函数记为 \(u_n\)。
两端值之差为 \(2^{-n}\)，故
\[
 u_n'=(3/2)^n\cdot\mathbf1_{C_n}\quad\text{几乎处处}.
\]
函数 \(u_n\) 连续且分段线性，因而属于 \(W^{1,1}\subseteq SBV\)。
它在每个基本区间内与 \(F_C\) 相差至多 \(2^{-n}\)，
在间隙上则相等，所以一致收敛成立。
又因为 \(m_1(C_n)=(2/3)^n\)，
\[
 \int_0^1|u_n'|\,\mathrm dx=1=|DF_C|\bigl((0,1)\bigr).
\]
这正是严格收敛，但 \(D^cF_C=\rho_C\ne0\)。
\end{proof}

这组函数没有跳跃，且梯度的 \(L^1\) 范数一致有界，
仍能把全部导数集中成奇异连续测度。
若 \(p>1\)，同一计算给出
\[
 \int_0^1|u_n'|^p\,\mathrm dx=(3/2)^{n(p-1)}\longrightarrow\infty.
\]
超线性的梯度控制正是后面紧性条件中的一个实质要求。

\subsection{紧性}
\label{b1:subsec:280403}

\begin{theorem}[Ambrosio 紧性：有界值域与超线性梯度（深层紧性定理）]
\label{b1:thm:ambrosio-sbv-compactness}
设 \(\Omega\) 为有界开集，\(1<p<\infty\)，\(u_j\in SBV(\Omega)\)，且
\[
 \sup_j\left\{\|u_j\|_{L^\infty(\Omega)}
       +\int_\Omega|\nabla u_j|^p\,\mathrm dx
       +\mathcal H^{d-1}(J_{u_j})\right\}<\infty.
\]
则存在子列与 \(u\in SBV(\Omega)\cap L^\infty(\Omega)\)，使
\[
 u_j\longrightarrow u\quad\text{在每个 }L^r(\Omega),\ 1\le r<\infty,
 \quad
 \nabla u_j\rightharpoonup\nabla u\quad\text{在 }L^p(\Omega;\mathbb R^d).
\]
此外，\(D^au_j\) 与 \(D^ju_j\) 分别\WeakStar{}收敛到 \(D^au\) 与 \(D^ju\)，
测度收敛均以 \(C_0(\Omega;\mathbb R^d)\) 为测试类。
\end{theorem}
\term[ambrosio-sbv-jinxing]{Ambrosio 紧性定理}[Ambrosio compactness theorem]
比一般有界变差函数的紧性多给出一个结论：极限仍没有 Cantor 部分。
\begin{proofsketch}
以下展开这个特殊版本的闭合步骤；所依赖的额外深层工具是
\cref{b1:thm:volpert-chain-rule}，不把它当作已经由普通 Sobolev 链律推出的结果。
记统一的值域界为 \(M\)，跳跃面积界为 \(S\)。
两侧近似值也在 \([-M,M]\) 中，故
\[
 |Du_j|(\Omega)
 \le m_d(\Omega)^{1-1/p}\|\nabla u_j\|_{L^p(\Omega)}
      +2M S.
\]
于是 \(BV\) 范数一致有界。由\cref{b1:thm:bv-local-compactness}取局部 \(L^1\) 收敛子列。
对于内部紧集穷竭 \(K_n\uparrow\Omega\)，
\(\int_{\Omega\setminus K_n}|u_j|\le M\,m_d(\Omega\setminus K_n)\) 一致趋零；
所以\cref{b1:cor:bv-global-compactness-with-tails}给出全域 \(L^1\) 收敛，
极限 \(u\in BV(\Omega)\) 且 \(|u|\le M\)。
再取子列，使 \(u_j\to u\) 几乎处处，且 \(\nabla u_j\rightharpoonup h\) 于 \(L^p\)。
尚需识别 \(h\) 并排除 \(D^cu\)。

取任意 \(\Phi\in C^1(\mathbb R)\)，满足 \(0\le\Phi\le1\) 且 \(\Phi'\) 有界。
由链律和 \(D^cu_j=0\)，测度
\[
 T_j^\Phi=D\bigl(\Phi(u_j)\bigr)-\Phi'(u_j)\nabla u_j\,m_d
\]
的总变差不超过 \(S\)。
函数 \(\Phi(u_j)\to\Phi(u)\) 于 \(L^1\)。
若 \(p'=p/(p-1)\)，对每个 \(q\in L^{p'}\)，
控制收敛给出
\(\Phi'(u_j)q\to\Phi'(u)q\) 于 \(L^{p'}\)，
因而 \(\Phi'(u_j)\nabla u_j\rightharpoonup\Phi'(u)h\)。
对紧支集光滑向量场取极限，再用变差的测试表述，得到
\[
 \left|D\bigl(\Phi(u)\bigr)-\Phi'(u)h\,m_d\right|(\Omega)\le S.
\]
再对 \(u\) 用链律，绝对连续、Cantor、跳跃三项互相奇异，因此
\[
 \int_\Omega|\Phi'(\widetilde u)|\,\mathrm d\sigma\le S,
 \quad
 \sigma=|\nabla u-h|\,m_d+|D^cu|.
\]
有限近似极限在 \(\sigma\)-几乎处处存在，故积分不依赖其余点的补值。
分别取
\[
 \Phi_N(t)=\frac{1+\sin(Nt)}2,\quad
 \Psi_N(t)=\frac{1+\cos(Nt)}2.
\]
因为 \(|\Phi_N'(t)|+|\Psi_N'(t)|\ge N/2\)，两次估计相加即得
\[
 \frac N2\,\sigma(\Omega)\le2S.
\]
令 \(N\to\infty\)，便有 \(\sigma=0\)，即 \(h=\nabla u\) 且 \(D^cu=0\)。
这个振荡测试正是所引论文 Propositions 2.3、2.5 在固定面积界下的核心机制。

最后，\(|u_j-u|\le2M\) 给出
\(\|u_j-u\|_r^r\le(2M)^{r-1}\|u_j-u\|_1\)，得到全部有限 \(r\) 的强收敛。
绝对连续部分的收敛来自梯度弱收敛，
跳跃部分的收敛来自 \(D^ju_j=Du_j-D^au_j\)
与\cref{b1:prop:bv-distributional-weak-star}。
\end{proofsketch}

这一定理为自由间断能量的直接法提供紧性，但尚未给出跳跃面积的下半连续性。
下一节先明确具体能量与容许类，再补上这项测度论接口。
